\chapter{Source Code}
\label{app:changed-source-code}

This appendix provides the key source code for the implemented optimisations. The code for the three core optimisations is from the combined (\texttt{all}) configuration variant, which includes all three optimisations. The code for the weak-field mechanism is from the \texttt{weak\_fields} configuration variant. Complete sources for all variants are available in the \texttt{patches/} directory of the project repository.

\lstset{basicstyle=\ttfamily\footnotesize,columns=flexible}

\section{Dynamic-array Data Structure}
\label{app:dynamic-array-source}

\noindent\textbf{File:} \texttt{zWeakRefArray.hpp}

\lstinputlisting[
	nolol,
	language=C++
]{../patches/all/src/hotspot/share/gc/z/zWeakRefArray.hpp}

\section{Reference Processor}
\label{app:refproc-impl-source}

\subsection{Interface}
\label{app:refproc-interface}

\noindent\textbf{File:} \texttt{zReferenceProcessor.hpp}

\lstinputlisting[
	nolol,
	language=C++
]{../patches/all/src/hotspot/share/gc/z/zReferenceProcessor.hpp}

\subsection{Implementation Excerpts}
\label{app:refproc-excerpts}

The following excerpts from \texttt{zReferenceProcessor.cpp} show the constructor and per-worker storage setup, the discovery path with queue-less weak-reference routing, the processing paths, and the queue sentinel initialisation.

\medskip
\noindent\textbf{Constructor and per-worker storage:}

\lstinputlisting[
	nolol,
	language=C++,
	firstline=132,
	lastline=170
]{../patches/all/src/hotspot/share/gc/z/zReferenceProcessor.cpp}

\noindent\textbf{Discovery path with queue-less routing:}

\lstinputlisting[
	nolol,
	language=C++,
	firstline=288,
	lastline=370
]{../patches/all/src/hotspot/share/gc/z/zReferenceProcessor.cpp}

\noindent\textbf{Processing paths for separate discovered lists and dynamic arrays:}

\lstinputlisting[
	nolol,
	language=C++,
	firstline=419,
	lastline=475
]{../patches/all/src/hotspot/share/gc/z/zReferenceProcessor.cpp}

\noindent\textbf{Queue sentinel initialisation and queue checks:}

\lstinputlisting[
	nolol,
	language=C++,
	firstline=680,
	lastline=713
]{../patches/all/src/hotspot/share/gc/z/zReferenceProcessor.cpp}

\section{Weak Fields}
\label{app:weak-fields}

\subsection{Annotation}
\label{app:weak-fields-annotation}

\noindent\textbf{File:} \texttt{java/lang/ref/weak.java}

\lstinputlisting[
	nolol,
	language=Java
]{../patches/weak_fields/src/java.base/share/classes/java/lang/ref/weak.java}

\subsection{Data Structure}
\label{app:weak-fields-array}

\noindent\textbf{File:} \texttt{zWeakFieldArray.hpp}

\lstinputlisting[
	nolol,
	language=C++
]{../patches/weak_fields/src/hotspot/share/gc/z/zWeakFieldArray.hpp}

\subsection{Reference Processor}
\label{app:weak-fields-refproc}

\noindent\textbf{File:} \texttt{zReferenceProcessor.hpp}

\lstinputlisting[
	nolol,
	language=C++
]{../patches/weak_fields/src/hotspot/share/gc/z/zReferenceProcessor.hpp}

\medskip
The following excerpts from \texttt{zReferenceProcessor.cpp} show the discovery filtering, the discovery function, and the processing function.

\medskip
\noindent\textbf{Discovery filtering:}

\lstinputlisting[
	nolol,
	language=C++,
	firstline=212,
	lastline=226
]{../patches/weak_fields/src/hotspot/share/gc/z/zReferenceProcessor.cpp}

\noindent\textbf{Weak-field discovery:}

\lstinputlisting[
	nolol,
	language=C++,
	firstline=326,
	lastline=360
]{../patches/weak_fields/src/hotspot/share/gc/z/zReferenceProcessor.cpp}

\noindent\textbf{Weak-field processing:}

\lstinputlisting[
	nolol,
	language=C++,
	firstline=393,
	lastline=433
]{../patches/weak_fields/src/hotspot/share/gc/z/zReferenceProcessor.cpp}

\subsection{Marking Closure}
\label{app:weak-fields-marking}

The following excerpt from \texttt{zMark.cpp} shows how weak-annotated fields are detected and diverted to the weak-field discovery path during object marking.

\medskip
\noindent\textbf{Field annotation check and discovery diversion:}

\lstinputlisting[
	nolol,
	language=C++,
	firstline=297,
	lastline=362
]{../patches/weak_fields/src/hotspot/share/gc/z/zMark.cpp}
